\section{Động lực và bối cảnh}
\begin{frame}{Từ sequential recommendation đến MLLM-MSR}
  \begin{columns}[T,onlytextwidth]
    \begin{column}{0.31\textwidth}\centering
      \textbf{Sequential recommendation}\par\vspace{0.2cm}
      \small Lịch sử tương tác có thứ tự: click, mua hoặc xem. Mục tiêu là dự đoán item tiếp theo dựa trên pattern theo thời gian.
    \end{column}
    \begin{column}{0.31\textwidth}\centering
      \textbf{Multimodal recommendation}\par\vspace{0.2cm}
      \small Một item không chỉ có ID/title mà còn có ảnh bìa, product image hoặc video cover.
    \end{column}
    \begin{column}{0.31\textwidth}\centering
      \textbf{MLLM}\par\vspace{0.2cm}
      \small Có khả năng nối image và language vào cùng không gian semantic — nhưng chưa tự nhiên hiểu dynamics của lịch sử dài.
    \end{column}
  \end{columns}
  \vfill\centering\Large\textcolor{Logo1}{Câu hỏi: làm sao để MLLM vừa nhìn thấy item, vừa hiểu người dùng thay đổi theo thời gian?}
\end{frame}

\begin{frame}{Ba khoảng trống mà paper nhắm tới}
  \begin{enumerate}
    \item \textbf{Nhiều ảnh theo thứ tự:} đưa cả chuỗi image vào MLLM vừa tốn context vừa khó ghép đúng ảnh với item.
    \item \textbf{Preference biến đổi:} user có thể chuyển từ một chủ đề/visual style sang chủ đề khác; một summary tĩnh dễ lỗi thời.
    \item \textbf{MLLM chưa phải recommender:} model pretrained hiểu ảnh và chữ, nhưng chưa được tối ưu cho xác suất tương tác với candidate item.
  \end{enumerate}
  \vspace{0.25cm}\begin{block}{Trực giác của MLLM-MSR}
    Nén mỗi item thành mô tả đa phương thức bằng ngôn ngữ, cập nhật user preference theo kiểu recurrent, rồi fine-tune một MLLM để dự đoán yes/no.
  \end{block}
\end{frame}

